\section{Earlier Work}
\label{sec:related}

\subsection{Transposing a matrix in place}

Catanzaro, Keller and Garland~\cite{catanzaro2014decomposition} split a matrix
transpose into row and column shuffles, with a rotation in between. Their
method needs a temporary buffer. Its cycle-walking step runs one block per
cycle. So it can only spread as wide as the number of cycles. They say this
themselves.

Gustavson, Karlsson and K\aa{}gstr\"{o}m~\cite{gustavson2012cycleleader} work
out where each cycle starts using number theory. They do not have to search for
it. We use the same idea for square blocks. There a transpose is its own
inverse, so the test becomes one comparison.

Sung et al.~\cite{sung2012pttwac} and G\'{o}mez-Luna et
al.~\cite{gomezluna2016inplace} pad dimensions to sizes that factor well.
Padding makes the array bigger. That works against saving memory. In our test
the worst case asked for many times the size of the tensor.

\subsection{Permuting more than two dimensions in place}

The closest work is ITTPD~\cite{cheng2025ittpd}. It breaks an $N$-D permutation
into block transpositions called $\rho(i,j,k)$. \textbf{That is the same
operation as our box swap. We do not claim it as ours.}

We differ in two ways.

First, ITTPD runs each step with Catanzaro's two-dimensional method. That needs
a temporary buffer of $\max(M,K)\cdot W$. It also goes over the data at least
once per step. We walk cycles directly. By Section~\ref{sec:fixedpoints} we can
go over only part of the data.

Second, ITTPD must choose its breakdown to fit a memory budget. When two
choices tie, it picks on memory rather than on time. Our memory cost is tiny
whichever breakdown we pick. So we are free to pick the cheapest one.

Where a back group of axes survives, we were faster on \resWinIttpd{} of
\resWinIttpdOf{} cases. Where it does not, ITTPD is faster than we are.
Section~\ref{sec:limits} has those numbers.

\subsection{Copying out of place}

The thing we are replacing is PyTorch's own
\texttt{.contiguous()}~\cite{pytorch_contiguous_2016,pytorch_2x}. It is fast. It
spreads perfectly over a GPU. It pays for that with a second tensor. Work on
memory pressure in large-model training~\cite{rajbhandari2020zero} and the
PyTorch communication document~\cite{pytorch_rfc_comms_2026} both treat that
cost as a real problem, not a theoretical one.
